% =====================================================================
% SEÇÃO DE METODOLOGIA  (cole dentro do seu documento principal)
%
% Pacotes necessários no preâmbulo:
%   \usepackage[utf8]{inputenc}
%   \usepackage[T1]{fontenc}
%   \usepackage[brazil]{babel}
%   \usepackage{amsmath, amssymb}
% =====================================================================

\chapter{Metodologia}

\section{Base de dados}

\begingroup\sloppy
Utilizou-se uma base de consumo de energia elétrica residencial com
$n = 17.179$ domicílios e 16 variáveis numéricas. A variável resposta é o
consumo em quilowatt-hora (\texttt{kwh}) e as 15 variáveis preditoras se
dividem em quatro grupos: \emph{moradores} (\texttt{moradores});
\emph{clima} (\texttt{graus\_dia\_frio\_f}, \texttt{graus\_dia\_calor\_f},
\texttt{temp\_inverno\_dia\_c}, \texttt{temp\_inverno\_noite\_c});
\emph{características do imóvel} (\texttt{total\_comodos},
\texttt{area\_casa\_m2}); e \emph{quantidade de aparelhos}
(\texttt{ventiladores\_teto}, \texttt{tvs}, \texttt{geladeiras},
\texttt{freezers}, \texttt{microondas}, \texttt{lava\_louca},
\texttt{maquina\_lavar}, \texttt{secadora}). A base foi recebida já
pré-processada e verificou-se que não há valores ausentes nem registros
duplicados. % TODO: citar aqui a fonte original dos dados.
\par\endgroup

\section{Regressão linear múltipla}

Como o objetivo é prever uma grandeza contínua (o consumo em kWh) a partir de
características do domicílio, adotou-se a regressão linear múltipla, por ser
simples e interpretável: cada coeficiente indica quanto o consumo varia ao se
alterar uma variável, mantidas as demais constantes. Uma regressão simples
seria insuficiente, pois a variável mais correlacionada com o consumo (a área
da casa, com $r \approx 0{,}39$) explicaria no máximo $r^2 \approx 15\%$ da
variância; as demais variáveis têm correlações moderadas e complementares. O
modelo é
\begin{equation}
  y_i = \beta_0 + \beta_1 x_{i1} + \dots + \beta_p x_{ip} + \varepsilon_i,
  \qquad i = 1, \dots, n,
  \label{eq:modelo}
\end{equation}
em que $y_i$ é o consumo do domicílio $i$, $x_{ij}$ é o valor da $j$-ésima
preditora, $\beta_0$ é o intercepto, $\beta_j$ são os coeficientes,
$\varepsilon_i$ é o erro aleatório e $p = 15$.

\paragraph{Por que um hiperplano.} Com $p = 1$ preditora, a
Equação~\eqref{eq:modelo} descreve uma reta no plano $(x, y)$; com $p = 2$,
um plano no espaço $(x_1, x_2, y)$. Em geral, os pontos
$(x_1, \dots, x_p, \hat{y})$ que satisfazem uma equação linear nas $p$
preditoras formam um subespaço de dimensão $p$ dentro de um espaço de
dimensão $p + 1$, ou seja, com uma dimensão a menos que o espaço que o contém.
Isso é, por definição, um hiperplano. Aqui os dados vivem em um espaço de
$16$ dimensões (15 preditoras mais o consumo) e a superfície ajustada é um
hiperplano de dimensão $15$; a reta da regressão simples é o caso particular
$p = 1$. Por isso a equação obtida não pode ser desenhada como uma única reta
em um gráfico bidimensional.

\section{Estimação e treinamento}

Os coeficientes são estimados pelo método dos mínimos quadrados ordinários,
que minimiza a soma dos quadrados dos resíduos:
\begin{equation}
  \hat{\boldsymbol{\beta}} = \arg\min_{\boldsymbol{\beta}}
  \sum_{i=1}^{n} (y_i - \hat{y}_i)^2
  = (\mathbf{X}^{\top}\mathbf{X})^{-1}\mathbf{X}^{\top}\mathbf{y},
  \label{eq:mqo}
\end{equation}
em que $\mathbf{X}$ é a matriz de preditoras (com uma coluna de uns para o
intercepto) e $\mathbf{y}$ o vetor de consumos. A implementação foi feita em
Python com a classe \texttt{LinearRegression} do \texttt{scikit-learn}, que
resolve o problema por decomposição em valores singulares, numericamente mais
estável que inverter $\mathbf{X}^{\top}\mathbf{X}$.

Os dados foram divididos aleatoriamente em 80\% para treino (13.743
domicílios) e 20\% para teste (3.436), com semente fixa
(\texttt{random\_state = 42}) para garantir reprodutibilidade. O modelo foi
ajustado apenas com o conjunto de treino; o de teste serviu para medir o
desempenho em dados não vistos. Não houve padronização das variáveis, pois
ela não altera as previsões da regressão por mínimos quadrados e, sem ela, os
coeficientes permanecem nas unidades originais (por exemplo, kWh por
morador). Além disso, aplicou-se validação cruzada com 5 partições sobre a
base completa, reportando-se a média e o desvio padrão do $R^2$.

\section{Métricas de avaliação}

Sejam $y_i$ o valor observado, $\hat{y}_i$ o previsto, $\bar{y}$ a média dos
valores observados, $n$ o número de amostras avaliadas e $p$ o número de
preditoras. As métricas foram calculadas nos conjuntos de treino e de teste:

\begin{description}
  \item[MSE (erro quadrático médio)] $= \frac{1}{n}\sum_{i=1}^{n}(y_i - \hat{y}_i)^2$.
  Média dos quadrados dos erros, em kWh\textsuperscript{2}. Penaliza muito
  mais os erros grandes e é a quantidade que o método dos mínimos quadrados
  minimiza.

  \item[RMSE (raiz do erro quadrático médio)] $= \sqrt{\mathrm{MSE}}$.
  Volta à unidade original (kWh) e pode ser lida como o erro típico do
  modelo, com maior peso para os erros grandes.

  \item[MAE (erro absoluto médio)] $= \frac{1}{n}\sum_{i=1}^{n}\lvert y_i - \hat{y}_i\rvert$.
  Quanto o modelo erra, em kWh, em média, para mais ou para menos. É menos
  sensível a valores extremos que o RMSE; um RMSE muito maior que o MAE indica
  a presença de alguns erros grandes.

  \item[$R^2$ (coeficiente de determinação)] $= 1 - \dfrac{\sum_{i}(y_i - \hat{y}_i)^2}{\sum_{i}(y_i - \bar{y})^2}$.
  Proporção da variância do consumo explicada pelo modelo, comparando seu erro
  com o de um modelo ingênuo que sempre prevê a média: $R^2 = 1$ indica
  previsão perfeita, $R^2 = 0$ indica desempenho igual ao de prever a média, e
  valores negativos (possíveis no teste) indicam desempenho pior que isso.

  \item[$R^2$ ajustado] $= 1 - (1 - R^2)\dfrac{n - 1}{n - p - 1}$.
  Desconta o número de preditoras. O $R^2$ comum nunca diminui ao se adicionar
  uma variável, mesmo irrelevante; o ajustado só aumenta se a variável melhora
  o modelo mais do que o esperado ao acaso.

  \item[$R^2$ da validação cruzada] Média e desvio padrão do $R^2$ nas 5
  partições. Uma média próxima do $R^2$ de teste e um desvio pequeno indicam
  que o desempenho não depende de uma divisão particular dos dados.
\end{description}

Comparar as métricas de treino e de teste também permite detectar
sobreajuste: se o desempenho no treino for muito superior ao do teste, o
modelo está memorizando os dados em vez de generalizar.